% Philosophy of Science review-copy adaptation.
% Venue-neutral source remains paper/main.tex.
% Abstract is <=100 words; author metadata intentionally omitted for anonymous review.
\documentclass[12pt]{article}

\usepackage[top=4cm,bottom=4cm,left=3cm,right=2.5cm]{geometry}
\usepackage[T1]{fontenc}
\usepackage[utf8]{inputenc}
\usepackage{lmodern}
\usepackage{microtype}
\usepackage{amsmath,amssymb}
\usepackage{enumitem}
\usepackage{setspace}
\usepackage{footmisc}
\usepackage{natbib}
\usepackage{xurl}
\usepackage[hidelinks]{hyperref}

\title{When Is a Formal Difference Admissible Evidence for an Individuation Claim?\\
\large Representation, Evidential Bridges, and Restricted Tests}
\author{}
\date{}

\newcommand{\Indist}{\mathrel{\equiv}}
\newcommand{\Dist}{\mathrel{\#}}
\doublespacing
\renewcommand{\footnotelayout}{\doublespacing\normalsize}

\begin{document}
\raggedright
\maketitle

\begin{abstract}
A formal difference between representations is not by itself evidence that represented targets are distinct. We formulate an individuation-specific admissibility framework separating representation features, target reference, evidential bridges, and admitted tests. A representation-level feature may contribute only when an independently justified experimental, causal, measurement, or semantic bridge connects it to the target; declaration alone does not supply warrant. A UMI sequencing case contrasts arbitrary read identifiers with experimentally grounded molecular tags. A small Lean artifact audits bridge dependence, relabeling invariance, and test-relative refinement. The framework states necessary, not sufficient, conditions for operational individuation claims.
\end{abstract}
\newpage

\section{Introduction}

A formal model can contain two different representations without thereby providing evidence for two different target-level individuals. Variables, constructors, identifiers, labels, and carrier elements can differ because of how a model is encoded rather than because the represented world contains two individuals.

The problem is evidential: when may a difference inside a representation be used as evidence for an operational individuation claim about the represented target?

Representation-level inequality alone is insufficient. A representation-level feature can contribute only when a defensible evidential bridge connects that feature to the target, and the resulting target-relevant difference is detectable under the tests admitted by the analysis.

The condition is necessary, not sufficient. A target-linked feature can still be unreliable, confounded, noisy, or badly calibrated. The framework asks a prior dependency question: has the argument shown why this formal difference is evidence about the target at all?

This yields a practical audit:
\begin{quote}
\textbf{If a representation-level identifier were arbitrarily reassigned while all independently justified target-relevant relations were held fixed, should the individuation conclusion change?}
\end{quote}
If the answer is yes, the argument owes an account of the evidential bridge from that identifier to the target.

The closest conceptual comparison is target-directed work on scientific representation and theoretical equivalence \citep{Suarez2004Inferential,Contessa2007Surrogative,Nguyen2017RepresentationEquivalence}. The present paper asks a narrower downstream question: when a difference in the representation itself is offered as evidence for individuation, what must connect it to the target before the inference is admissible?

\subsection{Contributions}

The paper contributes an explicit evidential-bridge condition, a relabeling-invariance diagnostic illustrated by UMI-based molecular counting, and a small mechanized dependency audit. The mathematical content is intentionally elementary.

\section{Formal Setting}

\subsection{Representations and targets}

Let \(P\) be a set of representations, \(T\) a target domain, and
\[
r:P\to T
\]
a representation-to-target map. Distinct elements of \(P\) are not assumed to represent distinct elements of \(T\).

The map \(r\) is a semantic input. The formal core is restricted to one target value per representation and does not infer \(r\) from observations.

\subsection{Representation features and evidential bridges}

Let
\[
f:P\to V
\]
be a representation-level feature and
\[
\phi:T\to V
\]
a target-level property. In the exact setting, define
\[
W(f,\phi)
\quad\Longleftrightarrow\quad
\forall p\in P,\;f(p)=\phi(r(p)).
\]

Then
\[
W(f,\phi)\land f(a)\neq f(b)
\quad\Longrightarrow\quad
r(a)\neq r(b).
\]

The implication is elementary. Its value is structural: the bridge makes explicit what must connect a presentation-level difference to the target. Declaration does not create epistemic warrant. Accepting such a bridge in a scientific application requires an independent experimental, causal, measurement, calibration, or semantic justification.

\subsection{Admitted tests}

Let \(Q\) be a class of target-sensitive tests and \(O(q,t)\) the outcome of test \(q\) on target \(t\). Let \(A\subseteq Q\) be a family admitted by an analysis. Admission itself requires an independently defensible connection between the test and the target claim.

Define
\[
a\Indist_A b
\Longleftrightarrow
\forall q\in A,\;O(q,r(a))=O(q,r(b)),
\]
and
\[
a\Dist_A b
\Longleftrightarrow
\exists q\in A:\;O(q,r(a))\neq O(q,r(b)).
\]

These relations are operational and test-relative, not definitions of metaphysical numerical identity.

\section{Scientific Case: UMI-Based Molecular Counting}

PCR amplification can produce many downstream sequencing reads from one pre-amplification DNA or RNA molecule. A hypothetical data set containing 100 read records therefore does not, merely by containing 100 records, support the conclusion that 100 source molecules were present.

\citet{Kivioja2012UMI} introduced unique molecular identifiers (UMIs) to support molecule counting by attaching molecular tags before amplification. \citet{SmithHegerSudbery2017UMITools} show why the tags still require error-aware analysis: UMI sequences can themselves contain sequencing errors, and naive token equality can misidentify PCR duplicates.

An ordinary software read identifier distinguishes records but supplies no bridge to source-molecule plurality. A pre-amplification UMI differs because the experimental procedure is designed to connect the tag to a source molecule before PCR generates downstream multiplicity.

In the idealized model, \(P\) is the set of read records, \(T\) the set of pre-amplification molecules, and \(r(p)\) the source molecule of read \(p\). Actual investigators do not directly observe \(r\); source assignment is inferred through protocol, alignment, barcode statistics, and error assumptions. The framework therefore audits a proposed source-assignment and evidential model rather than reconstructing deduplication from raw data.

Target relevance is not infallibility. Barcode collisions, sequencing errors, and other failures mean that same-UMI and different-UMI observations cannot be treated as unrestricted numerical-identity rules. Experimental provenance makes the tag evidentially relevant; the error model determines how particular observations should be interpreted.

\section{Test-Relative Distinguishability}

Let \(A_c\subseteq A_f\). Under fixed exact outcome semantics,
\[
a\Indist_{A_f}b\Longrightarrow a\Indist_{A_c}b.
\]
The converse need not hold. A restricted family can leave a pair unresolved that a richer justified family separates.

This is not a theorem about noisy evidence, statistical updating, model revision, or measurement disturbance.

\section{Mechanized Dependency Audit}

The Lean artifact is a dependency audit, not a claim of deep mathematical novelty.

The bridge layer checks that equal grounding forces equality of a bridged feature (\texttt{R16}) and that a difference in a bridged feature entails a grounding difference (\texttt{R17}). A deliberately representation-sensitive encoding probe admits no bridge because it separates two same-grounding presentations (\texttt{R18}), while a positive-control ground-tracking feature does admit one (\texttt{R19}).

The older checks remain supporting controls: \texttt{R1}--\texttt{R4} separate representation-sensitive discrimination from same-reference target-sensitive invariance; \texttt{R5}--\texttt{R7} cover test-family refinement; \texttt{R8}--\texttt{R10} cover the target-sensitive positive control; and \texttt{R11}--\texttt{R15} cover token and access-metadata invariance.

The result set is intentionally elementary. It makes hidden dependencies inspectable; it does not replace substantive scientific justification for accepting a bridge.

\section{Relation to Existing Work}

\subsection{Scientific representation and theoretical equivalence}

\citet{Suarez2004Inferential} and \citet{Contessa2007Surrogative} already treat scientific representation as target-directed and inferential. \citet{Nguyen2017RepresentationEquivalence} is the closest comparison: models can be compared by whether they license the same claims about the same targets.

The present paper claims no novelty for target-directed inference or same-target claim licensing. Its narrower question begins when a formal difference inside a representation is itself offered as evidence for target plurality. The bridge condition isolates that dependency: the feature must have a justified route to a target-level property rather than merely be formally discriminating.

\subsection{Individuation in scientific practice}

Practice-oriented work asks how scientists count, track, separate, manipulate, and present entities in concrete settings \citep{BuenoChenFagan2018,Waters2018AskNot,Love2018ExperimentalPractice}. \citet{Chen2018ExperimentalIndividuation} distinguishes ontological and epistemological modes of experimental individuation and explicitly discusses presentation.

The present contribution is narrower: when a representational device is used in an individuation practice, what justifies treating differences in that device as evidence about the target?

\subsection{Technical neighbors}

Representation independence \citep{Mitchell1986RepresentationIndependence}, observational and behavioral equivalence \citep{HennessyMilner1985,Rutten2000UniversalCoalgebra}, and work on identity and discernibility \citep{LadymanLinneboPettigrew2012,DieksVersteegh2008} constrain the interpretation of the present results. No new general equivalence theory or solution to numerical identity is claimed.

The negative result is evidential rather than eliminativist: representation-only or surplus structure may remain useful for other purposes \citep{NguyenTehWells2020Surplus}.

\section{Scope and Limitations}

The framework deliberately leaves several problems open.

The target domain and reference map are inputs; uncertain source assignment must be established by methods outside the exact formal core. The bridge condition is exact and deterministic, whereas real scientific warrants can be probabilistic and error-prone. Target relevance and test discrimination are necessary but not sufficient for epistemic warrant: reliability, calibration, confounding, background assumptions, and statistical decision rules remain substantive questions.

The formal core uses a function \(r:P\to T\); many-to-many, distributed, and probabilistic representation-target relations are outside the current scope. Observation and intervention are not equated. Operational distinguishability is not metaphysical numerical identity, and synchronic discrimination does not establish diachronic persistence.

\section{Relabeling Invariance as a Practical Audit}

For an identifier-like feature used in an individuation argument, ask whether arbitrary reassignment would alter the conclusion while independently justified target relations were held fixed. If it would, ask what experimental, causal, measurement, or semantic bridge explains why.

A representation-only discriminator changes bookkeeping or encoding while leaving the target-relevant model fixed. A bridged discriminator participates in a justified process connecting representation to target. It may therefore contribute evidence, subject to the reliability and error assumptions of that process.

The UMI case instantiates the contrast. Read identifiers fail the molecular relabeling audit. Pre-amplification UMIs are designed to pass it because their evidential role depends on physical tagging history rather than identifierhood.

\section{Formal Audit Summary}

For anonymous review, the mechanized claims are referred to only by neutral labels. The complete audit surface is \texttt{R1}, \texttt{R2}, \texttt{R3}, \texttt{R4}, \texttt{R5}, \texttt{R6}, \texttt{R7}, \texttt{R8}, \texttt{R9}, \texttt{R10}, \texttt{R11}, \texttt{R12}, \texttt{R13}, \texttt{R14}, \texttt{R15}, \texttt{R16}, \texttt{R17}, \texttt{R18}, and \texttt{R19}. Their definitions and scope are summarized in the anonymous artifact manifest.

\section{Conclusion}

A formal difference is not admissible evidence for an individuation claim merely because it distinguishes two representations.

The missing step is an evidential bridge. A representation-level feature may contribute to target-level individuation only when an independently justified experimental, causal, measurement, or semantic relation connects that feature to the target. A difference must then survive the relevant admitted tests. These are necessary conditions, not a complete epistemology of individuation.

UMI-based molecular counting shows why this matters. A read identifier and a molecular barcode can have the same formal type while playing very different evidential roles. The difference is explained by provenance and error-aware measurement practice, not by identifier syntax.

The formal contribution is correspondingly modest: make the bridge explicit, make relabeling invariance testable, and make hidden dependence on representation-level identity mechanically auditable.

% FUNDING_AND_DECLARATIONS: Before submission, add truthful Funding Statement
% and Declarations / competing-interests sections using author-supplied facts.
% Do not infer these values from repository or conversation context.
\section*{Acknowledgements}

AI tool disclosure: The author used OpenAI ChatGPT (GPT-5.6 Sol, accessed through ChatGPT on 22 September 2026) to assist with drafting, restructuring, translation, literature-search query formulation, and editorial revision. The author independently checked the arguments, references, formal claims, and final manuscript and takes full responsibility for the content.

\clearpage
\bibliographystyle{chicago}
\bibliography{references}

\end{document}
